\subsection{有限长度模}

回顾（I，§ 4，第 4 号，定义 7），一个 A-模 M 称为单模，如果它不约化为 0 且它不包含异于 M 和 \{0\} 的子模。一个 A-模 M 称为有限长度的，如果它有一个若尔当–赫尔德列 \((\mathbf{M}_t)_{0 \leqslant t \leqslant n}\) ，该列的商的个数 \(n\) （它不依赖于所考虑的 M 的若尔当–赫尔德列）称为 M 的长度（I，§4，第 7 号，定义 11）；我们将用 long(M) 或 long\_A(M) 表示它。一个约化为 0 的 A-模长度为 0；若 M 是一个非零的有限长度 A-模，则 long(M) \textgreater{} 0。

命题 16. 设 M 是一个 A-模，而 N 是 M 的一个子模；要使 M 具有有限长度，必须且只需 N 和 M/N 也具有有限长度，且此时

\[
\operatorname{long} (\mathbf {N}) + \operatorname{long} (\mathbf {M} / \mathbf {N}) = \operatorname{long} (\mathbf {M}).\tag{34}
\]

证明已在 I，§4，第 7 号，命题 10 中给出。

推论 1. 设 M 是一个有限长度的 A-模；对于 M 的子模 N，要使 N 等于 M，必须且只需 long(N) = long(M)。

推论 2. 设 \(u: \mathbf{M} \to \mathbf{N}\) 是一个 A-模同态。若 \(\mathbf{M}\) 或 \(\mathbf{N}\) 是有限长度的，则 \(\operatorname{Im}(u)\) 也是。若 \(\mathbf{M}\) 是有限长度的，则 \(\operatorname{Ker}(u)\) 也是，且

\[
\operatorname{long} (\operatorname{Im} (u)) + \operatorname{long} (\operatorname{Ker} (u)) = \operatorname{long} (\mathbf {M}).\tag{35}
\]

若 \(\mathbf{N}\) 是有限长度的，则 \(\operatorname{Coker}(u)\) 也是，且

\[
\operatorname{long} (\operatorname{Im} (u)) + \operatorname{long} (\operatorname{Coker} (u)) = \operatorname{long} (\mathbf {N}).\tag{36}
\]

推论 3. 设 \((\mathbf{M}_i)_{0\leqslant i\leqslant n}\) 是一个由有限长度的 A-模组成的有限族。若存在一个线性映射的正合列

\[
0 \longrightarrow \mathrm{M} _ {0} \stackrel {u _ {0}} {\longrightarrow} \mathrm{M} _ {1} \stackrel {u _ {1}} {\longrightarrow} \mathrm{M} _ {2} \longrightarrow \dots \longrightarrow \mathrm{M} _ {n - 1} \stackrel {u _ {n - 1}} {\longrightarrow} \mathrm{M} _ {n} \longrightarrow 0\tag{37}
\]

那么

\[
\sum_ {k = 0} ^ {n} (- 1) ^ {k} \operatorname{long} (\mathbf {M} _ {k}) = 0.\tag{38}
\]

该推论对 \(n = 1\) 是显然的，而对 \(n = 2\) 正是命题 16；我们对 \(n\) 作归纳。若 \(\mathbf{M}_{n-1}' = \operatorname{Im}(u_{n-2})\) ，则根据归纳假设，

\[
\sum_ {k = 0} ^ {n - 2} (- 1) ^ {k} \operatorname{long} (\mathbf {M} _ {k}) + (- 1) ^ {n - 1} \operatorname{long} (\mathbf {M} _ {n - 1} ^ {\prime}) = 0.
\]

另一方面，正合列 \(0 \to \mathbf{M}_{n-1}' \to \mathbf{M}_{n-1} \to \mathbf{M}_n \to 0\) 给出

\[
\operatorname{long} \left(\mathbf {M} _ {n - 1} ^ {\prime}\right) + \operatorname{long} \left(\mathbf {M} _ {n}\right) = \operatorname{long} \left(\mathbf {M} _ {n - 1}\right),
\]

由此得到关系式 (38)。

推论 4. 设 M 和 N 是 A-模 E 的两个有限长度子模；则 M + N 是有限长度的，且

\[
\operatorname{long} (\mathbf {M} + \mathbf {N}) + \operatorname{long} (\mathbf {M} \cap \mathbf {N}) = \operatorname{long} (\mathbf {M}) + \operatorname{long} (\mathbf {N}).\tag{39}
\]

只需对正合列 (29)（第 7 号）应用推论 3 即可。

\[
0 \rightarrow \mathbf {M} \cap \mathbf {N} \rightarrow \mathbf {M} \oplus \mathbf {N} \rightarrow \mathbf {M} + \mathbf {N} \rightarrow 0
\]

利用 \(\operatorname{long}(\mathbf{M} \oplus \mathbf{N}) = \operatorname{long}(\mathbf{M}) + \operatorname{long}(\mathbf{N})\) 这一事实（由 (34) 式）。

推论 5. 设 M 是一个 A-模，它是一个有限长度子模族 (N \(_{t}\) ) 的和。则 M 是有限长度的，且

\[
\operatorname{long} (\mathbf {M}) \leqslant \sum_ {i} \operatorname{long} \left(\mathrm{N} _ {i}\right).\tag{40}
\]

此外，要使 (40) 式两边相等，必须且只需 M 是诸 \(N_{t}\) 的直和。

已经看到（第 7 号，公式 (28)）存在一个典范满射线性映射 \(h \colon \bigoplus_{i} N_{i} \to M\) ；因此推论由 (35) 得出。

推论 6. 设 M 和 N 是 A-模 E 的两个子模，使得 E/M 和 E/N 是有限长度的模；则 E/(M ∩ N) 是有限长度的，且

\[
\text { long } (\mathbf {E} / (\mathbf {M} \cap \mathbf {N})) + \text { long } (\mathbf {E} / (\mathbf {M} + \mathbf {N})) = \text { long } (\mathbf {E} / \mathbf {M}) + \text { long } (\mathbf {E} / \mathbf {N}). \tag {41}
\]

只需对正合列 (30) 应用推论 3 即可。

\[
0 \rightarrow \mathrm{E} / (\mathrm{M} \cap \mathrm{N}) \rightarrow (\mathrm{E} / \mathrm{M}) \oplus (\mathrm{E} / \mathrm{N}) \rightarrow \mathrm{E} / (\mathrm{M} + \mathrm{N}) \rightarrow 0
\]

利用以下事实

\[
\operatorname{long} ((\mathbf {E} / \mathbf {M}) \oplus (\mathbf {E} / \mathbf {N})) = \operatorname{long} (\mathbf {E} / \mathbf {M}) + \operatorname{long} (\mathbf {E} / \mathbf {N}).
\]

推论 7. 设 \((\mathbf{M}_i)\) 是 A-模 \(\mathbf{E}\) 的子模的一个有限族，使得诸 \(\mathbf{E} / \mathbf{M}_i\) 都是有限长度的模。则 \(\mathbf{E} / (\bigcap_{i} \mathbf{M}_i)\) 是有限长度的，且

\[
\operatorname{long} \left(\mathrm{E} / \left(\bigcap_ {i} \mathrm{M} _ {i}\right)\right) \leqslant \sum_ {i} \operatorname{long} (\mathrm{E} / \mathrm{M} _ {i}).\tag{42}
\]

已知（公式 (27)）存在一个典范单射线性映射 \(\mathbf{E} / \left(\bigcap_{i}\mathbf{M}_{i}\right)\to \prod_{i}(\mathbf{E} / \mathbf{M}_{i})\)

注记。除第 7 号命题 9 外，第 2 至第 10 号的所有结果对任意带算子的交换群均成立，只需将陈述中的子模（相应地，商模）替换为稳定子群（相应地，由稳定子群得到的商群）；我们还约定将带算子的群的同态称为“线性映射”。第 7 号命题 9 的推论对带算子的交换群也成立：这对推论 4 和推论 5 是显然的，对推论 2 亦然，因为 \(\alpha\left(\sum_{i\in I}x_i\right)=\sum_{i\in I}\alpha x_i\) 对每个算子 \(\alpha\) 成立，而推论 3 由它得出。至于推论 1，只需注意若 N 是 F 的一个包含 \(u(S)\) 的稳定子群，则 \(\bar{u}^{-1}(N)\) 是 E 的一个包含 S 的稳定子群，因此 \(\bar{u}^{-1}(N)\) 包含 M，从而 \(u(M)\subset N\) 。

